\section{Design Analysis and Algorithm}
\label{sec:algorithm}

This section first recalls the G\&S~2012 baseline
(Section~\ref{sec:ks_baseline}), then examines the three design aspects
identified in the introduction (Section~\ref{sec:audit}).
The three modifications are presented in order of dependency:
column-generation LP (Section~\ref{sec:cg}),
feasibility-preserving stage transitions (Section~\ref{sec:proposition}),
and flat-$k$ arc selection (Section~\ref{sec:edge_selection}).
Fix~3 is presented before Fix~2 because Proposition~\ref{prop:feasibility}
is the structural reason the funnel's narrowing arc sets are safe to use:
the fill size $k$ can be reduced aggressively only because the incumbent
is always feasible by construction regardless of how small $k$ becomes.
Section~\ref{sec:full_algorithm} presents the complete FKS-CG-WS algorithm.

\subsection{Baseline Kernel Search (G\&S 2012)}
\label{sec:ks_baseline}

After solving the LP relaxation, facilities with positive LP dual
($\alpha^*_j > 0$) form the \emph{kernel} $\K = \Jone \cup \Jfrac$;
the remaining LP-closed facilities ($\alpha^*_j = 0$) form $\Jzero$ and are
organized into \emph{buckets}.
G\&S~2012 then proceeds as follows:

\begin{enumerate}
  \item \textbf{Kernel MILP.}
    Solve MILP restricted to facilities in kernel $\K$ and edges
    $\mathcal{E} = \{(j,i) : j \in \K,\; \rc{ji} \leq \gamma\}$,
    where $\gamma = \text{mean}\{\rc{ji} : \xLP_{ji} > 0\}$ is the mean
    reduced cost over LP-active arcs.
    An objective cutoff $z^H$ is imposed.
  \item \textbf{Bucket expansion.}
    Sort $\Jzero$ by $\rc{y_j}$ ascending.
    Add the next $l_{\text{buck}} = |\K|$ facilities as bucket $B_h$,
    connect to clients with $\rc{ji} \leq \gamma$, and solve
    MILP$(\K \cup B_h)$ with a forcing constraint $\sum_{j \in B_h} y_j \geq 1$.
    Repeat for up to $\bar{N}_B = 3$ buckets.
  \item \textbf{Kernel update.}
    If an improved solution opens facilities from the bucket, add them to the
    kernel.
    Facilities already in the kernel are removed after repeated non-selection,
    controlled by parameter $p$.
    An iterative variant restarts the solution phase when new facilities are
    added.
\end{enumerate}

This design gives several MILP-side controls over running time: the number of
buckets, their size, the removal rule, and the time limit per restricted MILP.

\subsection{Design Analysis: Three Refinement Opportunities}
\label{sec:audit}

Before presenting FKS, we examine three aspects of the G\&S~2012 design where
targeted modifications improve performance.
Each corresponds directly to one component of FKS.

\paragraph{LP construction cost.}
The full LP relaxation contains $|\J||\I|$ shipping variables.
At optimality, the large majority carry positive reduced cost and are
set to zero; they play no role in the LP solution and therefore provide
no useful signal to KS.
Constructing the full LP matrix is therefore unnecessary: column generation
recovers the same LP bound, facility partition, and reduced-cost rankings
without ever adding the irrelevant variables.
On the largest instances (e.g., $|\J|=2{,}000$, $|\I|=2{,}000$), the full LP
solve can dominate total running time, making KS slower than necessary before
the first MILP is even solved.

\paragraph{$\gamma$-threshold arc selection: size and bias.}
G\&S~2012 includes arc $(j,i)$ in a restricted MILP whenever its reduced cost
satisfies $\rc{ji} \leq \gamma$, where $\gamma$ is the mean reduced cost over
LP-active arcs.
Two concerns follow.
First, the number of arcs admitted by $\gamma$ depends on LP concentration and
can vary widely across instances and capacity regimes, giving unpredictable
subproblem sizes.
Second, consider a bucket facility $j \in \Jzero$ (LP dual $\alpha^*_j = 0$).
By equation~\eqref{eq:pricing}, its arc reduced cost is
$\rc{ji} = c_{ji} - \pi^*_i - \alpha^*_j = c_{ji} - \pi^*_i$: the capacity
dual contributes nothing.
For a kernel facility $j' \in \Jone$ with $\alpha^*_{j'} > 0$, the same
formula gives $\rc{j'i} = c_{j'i} - \pi^*_i - \alpha^*_{j'} < c_{j'i} - \pi^*_i$.
The capacity dual reduces kernel arc reduced costs, making them easier to pass
the $\gamma$ threshold.
Bucket arcs must clear a structurally higher hurdle with no capacity signal
subtracted, so fewer bucket arcs qualify for any given $\gamma$.
This treats bucket facilities less favorably than kernel facilities, giving
them a systematically smaller arc set in every restricted MILP.

\paragraph{Cold-start MILP transitions.}
G\&S~2012 passes no information between consecutive restricted MILPs.
When the arc set changes (between stages or between kernel and bucket MILPs),
the incumbent solution from the previous subproblem may not be feasible in the
next one, so the solver must find feasibility from scratch.
In the multi-stage setting this is wasteful: a good solution is available but
cannot be used because there is no guarantee it remains feasible after the
arc set is modified.

FKS addresses each aspect in turn through the components described in the
following sections.

\subsection{Column Generation for the LP}
\label{sec:cg}

\paragraph{Motivation.}
The full LP over all $|\J| \times |\I|$ shipping variables can dominate KS
runtime on large instances.
The MSCFLP LP solution is typically sparse: only a small fraction of shipping
variables carry positive flow at optimality.
Column generation exploits this structure by starting with a small column set and
adding only inactive variables that can improve the LP objective.

\paragraph{Column generation setup.}
We keep all $y_j$ variables in the model from the start (there are only
$|\J|$ of them).
Only the $x_{ji}$ shipping variables are generated on demand.

The \emph{restricted master problem} (RMP) at iteration $t$ contains:
\begin{itemize}
  \item All $y_j$ variables with their opening costs $f_j$.
  \item A subset $\mathcal{A}^t \subseteq \J \times \I$ of active
        shipping columns with costs $c_{ji}$.
  \item Demand constraints~\eqref{eq:demand}: $\sum_{j:(j,i)\in\mathcal{A}^t}
        x_{ji} = d_i$ for each $i$ (initially enforced via a small artificial
        variable to guarantee feasibility).
  \item Capacity constraints~\eqref{eq:capacity}: $\sum_{i:(j,i)\in\mathcal{A}^t}
        x_{ji} \leq q_j y_j$ for each $j$.
  \item VUB constraints~\eqref{eq:vub}: $x_{ji} \leq d_i y_j$ for each
        $(j,i) \in \mathcal{A}^t$.
\end{itemize}
When column $(j,i)$ is added to $\mathcal{A}^t$, its variable $x_{ji}$ enters
the demand constraint for client $i$, the capacity constraint for facility $j$,
and a new VUB constraint.

\paragraph{Initialization.}
We initialize $\mathcal{A}^0$ by selecting the $k_{\text{init}} = 30$
cheapest facilities (by shipping cost $c_{ji}$) for each client $i$.
This gives $30 \times |\I|$ initial columns (120{,}000 for a
$1{,}000 \times 4{,}000$ instance), which is typically sufficient to produce
a feasible RMP.

\paragraph{Pricing step.}
After solving the RMP at iteration $t$, let $\pi_i$ be the dual of demand
constraint $i$ and $\alpha_j$ the dual of capacity constraint $j$.
The reduced cost of any inactive shipping variable $x_{ji}$
($(j,i) \notin \mathcal{A}^t$) is:
\begin{equation}
  \rc{ji} = c_{ji} - \pi_i - \alpha_j.
  \label{eq:pricing}
\end{equation}
If $\rc{ji} < 0$, then adding $x_{ji}$ to the RMP can decrease the LP
objective.
We add \emph{all} inactive columns with $\rc{ji} < 0$ simultaneously, then
re-solve.
This is the ``full pricing'' variant of column generation.
It is more expensive per iteration than pricing a single column, but it avoids
long sequences of small RMP resolves and is effective when many columns have
negative reduced cost in early iterations.
When a shipping column is added, its VUB constraint is introduced at the same
time.
After convergence, reduced costs used for edge selection are taken from the
final LP solution: active generated columns use the solver's reduced costs,
while inactive columns use the final demand and capacity dual prices.
The computational study checks the CG LP against the full LP whenever the full
LP can be solved within the experimental protocol.

\paragraph{Termination.}
When no inactive column has $\rc{ji} < 0$ and all artificial variables are zero,
the current RMP solution is accepted as the LP solution passed to KS.
This follows from LP duality: at the optimal RMP dual $(\pi, \alpha)$,
every column with $\rc{ji} \geq 0$ can only increase the objective if added,
so no ungenerated shipping column is attractive under the pricing oracle.

\paragraph{Vectorized implementation.}
The pricing step~\eqref{eq:pricing} is a matrix subtraction.
Let $C \in \mathbb{R}^{|\J| \times |\I|}$ be the shipping cost matrix.
After solving the RMP, the reduced costs of all $|\J| \times |\I|$ pairs are:
\[
  \mathrm{RC} = C - \boldsymbol{\pi}^{\top}[\text{broadcasted}]
                  - \boldsymbol{\alpha}[\text{broadcasted}].
\]
All inactive entries of $\mathrm{RC}$ with value $< 0$ are added in one
operation.
For $|\J| = 1{,}000$ and $|\I| = 4{,}000$, this is a single NumPy operation on a
$1{,}000 \times 4{,}000$ matrix and is small relative to the RMP solves in our
target setting.

\paragraph{Convergence in practice.}
For each reported instance where the full LP is solved, we compare the CG
objective value, facility partition, and reduced-cost rankings with those
obtained from the full LP.
This check is important because the purpose of CG in FKS is acceleration, not a
change in the information used to guide Kernel Search.
Section~\ref{sec:study_lp} reports the resulting CG iteration counts, active
columns, LP times, and full-LP comparisons.

\paragraph{Remark on branch-and-price.}
Column generation applied to the LP relaxation does \emph{not} make FKS a
branch-and-price algorithm.
Branch-and-price is an exact method that solves an LP at every node of a
branch-and-bound tree~\citep{barnhart1998branch}.
FKS solves the LP exactly once (via CG), then runs heuristic MILPs with no
LP re-solve.
The heuristic character of FKS is unchanged; CG is purely a speed-up for
the single LP solve that KS already requires.

\subsection{Feasibility-Preserving Stage Transitions}
\label{sec:proposition}

When an arc set changes between stages, arcs used by the incumbent may be
absent from the new set, making the incumbent infeasible.
The solver must then rediscover feasibility from scratch.
We eliminate this by including all incumbent arcs in the next stage's edge
set by construction.

Let $(y^s, x^s)$ be the best solution found after Stage~$s$.
For Stage~$s{+}1$ with parameter $k_{s+1}$, we define the edge set as:
\begin{equation}
  \mathcal{E}^{\text{inc}}_{k_{s+1}}(x^s)
  = \bigcup_{i \in \I}
    \Bigl(
      \underbrace{\{(j, i) : x^s_{ji} > 0\}}_{\text{incumbent arcs}}
      \;\cup\;
      \underbrace{\text{top-}k_{s+1}\text{ by } \rc{ji}
      \text{ in } \K}_{\text{reduced-cost fill}}
    \Bigr).
  \label{eq:stage2_edges}
\end{equation}
Each later stage therefore contains the incumbent arcs plus a reduced-cost
fill set.
The edge count is at most $k_{s+1}|\I|$ plus the number of incumbent arcs
not already in the fill set.

\begin{proposition}[Feasibility Preservation]
\label{prop:feasibility}
Let $(y^s, x^s)$ be a feasible solution to the Stage~$s$ restricted MILP
with facility set $\mathcal{F}^s$ and edge set $\mathcal{E}^s$.
Let the Stage~$(s{+}1)$ facility set $\mathcal{F}^{s+1}$ contain every
facility with $y^s_j=1$.
Let $\mathcal{E}^{s+1} = \mathcal{E}^{\mathrm{inc}}_{k_{s+1}}(x^s)$ as in
equation~\eqref{eq:stage2_edges}.
Then $(y^s, x^s)$ is feasible in the Stage~$(s{+}1)$ restricted MILP.
\end{proposition}

\begin{proof}
By construction~\eqref{eq:stage2_edges}, for every client $i$ and every
facility $j$ with $x^s_{ji} > 0$, the arc $(j,i)$ belongs to
$\mathcal{E}^{s+1}$.
The assumption on $\mathcal{F}^{s+1}$ ensures every facility opened by
$y^s$ is available in the next restricted MILP.
We verify the three constraint types:
\begin{enumerate}
  \item \textbf{Demand.} $(y^s, x^s)$ satisfies $\sum_j x^s_{ji} = d_i$
    for every client $i$. The same equality holds in Stage~$(s{+}1)$
    because every arc used by $x^s$ is present.
  \item \textbf{Capacity.} $(y^s, x^s)$ satisfies
    $\sum_i x^s_{ji} \leq q_j y^s_j$. Since $y^s$ and $x^s$ are unchanged,
    this constraint is satisfied.
  \item \textbf{Variable upper bound.} $(y^s, x^s)$ satisfies
    $x^s_{ji} \leq d_i y^s_j$ for every arc with $x^s_{ji} > 0$.
    These arcs are present in $\mathcal{E}^{s+1}$, so the VUB constraints
    are present and satisfied.
\end{enumerate}
Therefore $(y^s, x^s)$ is feasible in Stage~$(s{+}1)$. \hfill$\square$
\end{proof}

\begin{corollary}
\label{cor:warmstart}
When $(y^s, x^s)$ is injected as a MIP start for Stage~$(s{+}1)$, the
solver receives a feasible incumbent by construction and can begin from a
known primal bound.
\end{corollary}

\begin{remark}
Proposition~\ref{prop:feasibility} does not depend on the LP solution
structure, the instance size, or the capacity ratio.
It holds for any multi-stage matheuristic that retains incumbent facilities
and builds edge sets as in equation~\eqref{eq:stage2_edges}, making the
feasibility-preserving construction a general design principle for
Kernel Search algorithms.
\end{remark}

\paragraph{Practical effect.}
Stage~1 uses \texttt{MIPFocus=1} (Gurobi's feasibility-first mode) because
no warm start is yet available.
Stages~2 and~3 inject the previous incumbent and use \texttt{MIPFocus=0}
(balanced), because a primal bound is already known.
This removes the need for later stages to rediscover feasibility after the
arc set has been narrowed.

The feasibility guarantee is also the structural reason the flat-$k$ funnel
in the next section is safe to use: because Proposition~\ref{prop:feasibility}
holds regardless of how narrow the fill set is, the fill size $k$ can be
reduced aggressively across stages without risking cold restarts.

\subsection{Edge Selection: The \texorpdfstring{Flat-$k$}{Flat-k} Criterion}
\label{sec:edge_selection}

Section~\ref{sec:proposition} establishes that all incumbent arcs must be
retained in the next stage's edge set.
This section specifies the remaining arcs: the reduced-cost fill in
equation~\eqref{eq:stage2_edges}.
We use a \emph{client-centric, flat-$k$} criterion: for each client $i$,
include the $k$ facilities with the lowest reduced cost $\rc{ji}$.

\paragraph{Stage~1 edge set.}
Stage~1 has no incumbent, so the edge set is a pure reduced-cost selection:
\begin{equation}
  \mathcal{E}^{\text{RC}}_{k_1}
  = \bigcup_{i \in \I}
    \bigl\{ (j, i) : j \text{ is among the } k_1 \text{ lowest-}\rc{ji}
    \text{ facilities in } \K \bigr\}.
  \label{eq:stage1_edges}
\end{equation}
We use $k_1 = 50$, giving 50{,}000 edges on a $1{,}000 \times 1{,}000$
instance and scaling linearly with $|\I|$ regardless of LP behavior.
Stages~2 and~3 use equation~\eqref{eq:stage2_edges} with $k_2=20$ and
$k_3=10$: the fill set shrinks while the incumbent arcs are always retained
by Proposition~\ref{prop:feasibility}.

\begin{proposition}[Guaranteed subproblem size]
\label{prop:size}
For any MSCFLP instance, any LP solution, and any $k \geq 1$, the
Stage~1 edge set satisfies $|\mathcal{E}^{\mathrm{RC}}_{k_1}| = k_1 \cdot |\I|$.
This bound is exact and holds regardless of LP concentration, capacity
ratio, or instance scale.
\end{proposition}

The proof is immediate: exactly $k_1$ arcs are selected per client and
clients are disjoint.
The Stage~1 subproblem size is fully controlled through $k_1$ alone.

\paragraph{Why flat-$k$ rather than a threshold $\gamma$?}
Section~\ref{sec:audit} gives two reasons.
First, the $\gamma$-threshold provides no guarantee on subproblem size.
Second, it is structurally biased: bucket facilities have $\alpha^*_j = 0$,
so their arc reduced costs receive no capacity-dual discount and face a
higher effective hurdle than kernel-facility arcs under the same $\gamma$.
The flat-$k$ criterion eliminates both problems: it selects exactly $k$ arcs
per client and ranks all facilities by the same reduced-cost ordering,
treating kernel and bucket facilities identically.

\paragraph{Why three stages?}
The specific values $k_1=50$, $k_2=20$, $k_3=10$ are a practical choice,
not a theoretically derived optimum.
What matters is that the sequence is decreasing (so the fill set narrows
across stages) and that the feasibility guarantee of
Proposition~\ref{prop:feasibility} holds at each transition regardless of
how small $k$ becomes.
A practitioner with more time or memory can raise all three values.

\subsection{Full Algorithm}
\label{sec:full_algorithm}

Algorithm~\ref{alg:fks} states FKS precisely.
All parameters not listed in the algorithm follow G\&S~2012 defaults.
We use $k_1 = 50$, $k_2 = 20$, $k_3 = 10$ and a 300-second time limit per
stage throughout.
When a bucket solution improves the incumbent, any opened bucket facility is
promoted into the kernel before the next stage; this keeps the incumbent's
facility set available for the feasibility-preserving transition.
In the pseudocode, $\mathcal{E}^{\text{inc}}_k(x;\mathcal{F})$ denotes the
same incumbent-preserving flat-$k$ construction as
equation~\eqref{eq:stage2_edges}, applied over facility set $\mathcal{F}$.

\begin{algorithm}[t]
\caption{Funnel Kernel Search (FKS)}
\label{alg:fks}
\begin{algorithmic}[1]
\Require Instance $(\J,\I,f,q,d,c)$; stage parameters $(k_s, T_s)_{s=1}^{S}$
\Ensure Feasible integer solution $(y^H, x^H)$

\State \textbf{LP phase:} Solve LP via column generation
       $\to (\xLP, \yLP, \bar{c})$; \quad $z^{\text{LP}} \gets$ LP objective
\State $\K \gets \Jone \cup \Jfrac$
\State Sort $\Jzero$ by $\rc{y_j}$ ascending $\to$ bucket list $L$
\State $z^H \gets +\infty$; \quad $x^{\text{prev}} \gets \emptyset$;
       \quad $y^H \gets \emptyset$

\For{$s = 1, 2, \ldots, S$}

  \State \textbf{Build edge set:}
  \If{$s = 1$}
    \State $\mathcal{E} \gets \mathcal{E}^{\text{RC}}_{k_1}$
           \hfill\Comment{top-$k_1$ by reduced cost per client}
    \State \texttt{MIPFocus} $\gets 1$
           \hfill\Comment{find any feasible incumbent quickly}
  \Else
    \State $\mathcal{E} \gets \mathcal{E}^{\text{inc}}_{k_s}(x^{\text{prev}})$
           \hfill\Comment{preserve incumbent arcs, fill to $k_s$}
    \State \texttt{MIPFocus} $\gets 0$
           \hfill\Comment{balanced; warm start already feasible}
  \EndIf

  \State \textbf{Kernel MILP:}
         $(z, y, x) \gets
         \text{MILP}(\K,\, \mathcal{E},\, z^H,\, \emptyset,\, T_s,\,
         \texttt{warm}=(y^H,x^H))$
  \If{$z < z^H$} \; $z^H, y^H, x^H, x^{\text{prev}} \gets z, y, x, x$
  \EndIf

  \If{$s > 1$} \hfill\Comment{bucket expansion (stages 2+ only)}
    \For{$h = 1, \ldots, \bar{N}_B$}
      \State $B_h \gets$ next $l_{\text{buck}} = |\K|$ facilities from $L$
      \If{$B_h = \emptyset$} \textbf{break} \EndIf
      \State $\mathcal{E}' \gets
             \mathcal{E}^{\text{inc}}_{k_1}(x^H;\,\K \cup B_h)$
             \hfill\Comment{flat-$k$ bucket coverage plus incumbent arcs}
      \State $(z, y, x) \gets
             \text{MILP}(\K \cup B_h,\, \mathcal{E}',\, z^H,\, B_h,\, T_s,\,
             \texttt{warm}=(y^H,x^H))$
      \If{$z < z^H$} \;
        $z^H, y^H, x^H, x^{\text{prev}} \gets z, y, x, x$;
        $\K \gets \K \cup \{j \in B_h : y_j = 1\}$
      \EndIf
    \EndFor
  \EndIf

\EndFor
\State \Return $(y^H, x^H)$
\end{algorithmic}
\end{algorithm}

\paragraph{Why three stages and these $k$ values?}
The specific values $k_1=50$, $k_2=20$, $k_3=10$ are a practical choice, not
a theoretically derived optimum.
The flat-$k$ design is intentionally parameterized: a practitioner with more
time or memory can raise all three values; a practitioner with tighter
constraints can lower them.
What matters structurally is only that the sequence is decreasing (so the
subproblem narrows progressively) and that incumbent arcs are preserved at each
transition (so feasibility is not broken).
The values used here were chosen to give a broad first-stage coverage on the
largest benchmark families ($k=50$ gives 100{,}000 arcs on a $1{,}000\times
2{,}000$ instance) while keeping later stages computationally tractable.
Because every stage after the first begins from a feasible warm start
(Proposition~\ref{prop:feasibility}), later stages spend their time
improving the incumbent rather than rediscovering feasibility.

\paragraph{Relation to G\&S~2012.}
FKS retains all structural components of G\&S~2012: the same kernel
definition, the same bucket expansion structure, and the same MILP
formulation.
The two changes are: (1)~the LP is solved via CG rather than by constructing the
full LP upfront, and (2)~the restricted-MILP phase is replaced by a three-stage
funnel with feasibility-preserving transitions.
The computational study evaluates these changes separately before reporting the
combined end-to-end effect.
